\chapter{Lead--Lag and Cross-Venue Trading}\label{ch:hf:lead-lag-and-cross-venue-trading}

When the index future moves, the fund follows within milliseconds and the constituent stocks a little later. Who leads depends on the hour, the
venue and the size of the move, and it has changed as the fastest firms have caught up. Between 2005 and 2011 the median life of an arbitrage between
the S\&P 500 future and its fund fell from 97 milliseconds to 7. In this chapter's simulated markets, where a fund follows its future half a second
late, the lead is recovered from the two mids to within 25 milliseconds. The future's information share is 97--98\%. A taker who trades the fund when
the future has moved and the fund has not earns 0.49 ticks a share if it acts within a millisecond. The edge is gone before the half-second is.

\section{Price discovery across markets}

Two prices of one thing (a future and the fund that tracks the same index, one stock on two venues, a coin on two exchanges) share a common
efficient price and differ by transitory noise. Price discovery is the process by which new information enters that common price; the market where it
enters first leads, and the other follows. Hasbrouck found, for US equity indices, that most price discovery happens in the E-mini futures for the S\&P
500 and the Nasdaq-100, that it is shared between the regular future and the fund for the S\&P 400, and that the S\&P 500 fund leads the sector funds
far more than they lead it. Mizrach and Neely found the Treasury futures carrying most of the price discovery of five- and ten-year notes from 1999 on.

At the highest frequencies the two prices do not even move together. Budish, Cramton and Shim measured, for the S\&P 500 future and its fund in 2011,
a median return correlation of 0.10 at 10 milliseconds and 0.008 at one millisecond, against nearly one over a minute: the Epps effect of One Quant
Book 4, chapter 21, turned into an opportunity for whoever reads the leader first.

\begin{definition}[Cross-venue lead]\label{def:hf:lead-lag-and-cross-venue-trading:lead}
A \emph{cross-venue lead}\index{cross-venue lead} is the time by which the prices of one market (a venue, or an instrument sharing the same
efficient price) move ahead of those of another: the lag at which the cross-correlation of their price changes peaks, estimated on their own
asynchronous ticks.
\end{definition}

\section{Information shares}

One Quant Book 10, chapter 8 defines the two measures of price discovery. The \emph{information share} of a market is the share of the variance of
the common efficient price's innovations that comes from that market's innovations; the \emph{component share} is that market's weight in the
common efficient price. Both come from a vector error-correction model (One Quant Book 4, chapter 20) of the two prices $p^1,p^2$ on a regular grid,
with the cointegrating vector fixed at $(1,-1)$:
\[
\Delta p_t=\alpha\,(p^1_{t-1}-p^2_{t-1})+\sum_{i=1}^{k}\Gamma_i\,\Delta p_{t-i}+e_t .
\]
The market that corrects towards the other has the large adjustment coefficient; the leader's is near zero.

\begin{proposition}[The two shares of a pair]\label{prop:hf:lead-lag-and-cross-venue-trading:shares}
With $\alpha=(\alpha_1,\alpha_2)$ and $\alpha_\perp=(\alpha_2,-\alpha_1)$, the component share of market 1 is $\gamma_1=\alpha_2/(\alpha_2-\alpha_1)$. With
$\Omega=\mathrm{Cov}(e_t)$ and $F$ its Cholesky factor for a given ordering of the markets, the information share of market 1 is
$[(\gamma^\top F)_1]^2/\gamma^\top\Omega\gamma$, and the two orderings give its lower and upper bounds.
\end{proposition}

\admitted
The derivation of the common-trend representation is Hasbrouck's and Gonzalo and Granger's; the form above uses the fact, for a pair with
cointegrating vector $(1,-1)$, that the common trend's weights are proportional to $\alpha_\perp$.

\Cref{fig:hf:lead-lag-and-cross-venue-trading:shares} shows the leader's shares on the chapter's markets. At a lead of 50 milliseconds, shorter
than the 0.1-second grid, the innovations of the two markets are simultaneous on the grid and the information share is not identified: its bounds are
0.52 and 0.73. At half a second they are 0.97 and 0.98; at one second both are 0.998. The bounds narrow as the grid resolves the lead, which is why
information shares on one-second data say little about a millisecond race.

\begin{omfigure}
\begin{tikzpicture}
\begin{axis}[width=11cm,height=5.2cm,xlabel={\small planted lead of A over B (seconds)},ylabel={\small share of A},xmode=log,log basis x=10,
  xtick={0.05,0.2,0.5,1},xticklabels={0.05,0.2,0.5,1},xmin=0.04,xmax=1.25,ymin=0.4,ymax=1.02,tick label style={font=\footnotesize},
  grid=major,grid style={gray!25},table/col sep=comma,legend style={font=\scriptsize,at={(0.5,-0.3)},anchor=north},legend columns=3]
\addplot[omDef,thick,mark=*,mark size=1.3pt] table[x=lag,y=is_high]{figdata/hft/08-lead-lag-and-cross-venue-trading/shares.csv};
\addplot[omDef,thick,dashed,mark=o,mark size=1.3pt] table[x=lag,y=is_low]{figdata/hft/08-lead-lag-and-cross-venue-trading/shares.csv};
\addplot[omThm,thick,mark=square*,mark size=1.3pt] table[x=lag,y=component]{figdata/hft/08-lead-lag-and-cross-venue-trading/shares.csv};
\legend{information share (upper),information share (lower),component share}
\end{axis}
\end{tikzpicture}

\omcaption{Price-discovery shares of the leading market A, from a vector error-correction model on a 0.1-second grid with ten lags, by the lead
planted in the simulation (log scale); mean of two twenty-minute sessions. Data: \texttt{hf\_leadlag.table}.}\label{fig:hf:lead-lag-and-cross-venue-trading:shares}
\end{omfigure}

\section{Estimating lag on asynchronous ticks}

Two markets' prices change at different moments; sampling both on a grid adds the Epps bias of One Quant Book 4, chapter 21. The lead--lag estimator
of One Quant Book 7, chapter 10 works on the ticks themselves: the Hayashi--Yoshida covariance of the two series' increments, with one series' clock
shifted by a candidate lag, normalised into a correlation; the lag with the largest correlation is the lead
(\cref{fig:hf:lead-lag-and-cross-venue-trading:ccf}).

\begin{omfigure}
\begin{tikzpicture}
\begin{axis}[width=11cm,height=5.2cm,xlabel={\small candidate lag of B behind A (seconds)},ylabel={\small cross-correlation},
  tick label style={font=\footnotesize},xmin=-1.5,xmax=1.5,ymax=0.4,grid=major,grid style={gray!25},table/col sep=comma]
\addplot[omDef,thick] table[x=lag,y=ccf]{figdata/hft/08-lead-lag-and-cross-venue-trading/ccf.csv};
\draw[omThm,dashed] (axis cs:0.5,0) -- (axis cs:0.5,0.4);
\node[font=\scriptsize,omThm,anchor=north east,fill=white,inner sep=1pt] at (axis cs:0.4,0.39) {planted lead 0.5 s};
\end{axis}
\end{tikzpicture}

\omcaption{Hayashi--Yoshida cross-correlation of the two mids' changes, on their own ticks, with B's clock moved back by each candidate lag, in
steps of 25 milliseconds; one twenty-minute session with a planted lead of half a second. Data: \texttt{hf\_leadlag.estimate}.}\label{fig:hf:lead-lag-and-cross-venue-trading:ccf}
\end{omfigure}

On two sessions for each planted lead, the estimates are 0.05 and 0.25 seconds for 0.05, 0.225 and 0.25 for 0.2, 0.475 and 0.575 for 0.5, and 1.075
and 1.1 for one second. A market can be too slow for this: in Book 7's default configuration each simulated book follows its efficient price over
several seconds, the peak of the cross-correlation is flat, and the same estimator returned leads of $-1.65$ and $+1.0$ seconds for a planted 50
milliseconds. The chapter's markets make informed traders four times as active and stale quotes cancel faster, so that each mid follows its efficient
price within a fraction of a second; only then is a sub-second lead identified.

\section{Trading the lead: quoting and taking}

A market maker in the follower uses the lead to protect its quotes (chapter 7's futures-led skew). A taker uses it to cross the follower's spread
before the follower moves. The taker's rule here: when the leader's mid has moved by at least a tick over the last second more than the follower's,
cross the follower's spread for one lot, and mark the trade at the follower's mid five seconds later, after half the spread and a fee of 0.1 tick.
Acting $d$ seconds after the leader's move (\cref{fig:hf:lead-lag-and-cross-venue-trading:edge}; planted lead half a second, 135 trades in two
sessions):

\begin{omfigure}
\begin{tikzpicture}
\begin{axis}[width=11cm,height=5.2cm,xlabel={\small taker's latency after the leader's move (milliseconds)},ylabel={\small edge (ticks a share)},
  xmode=log,log basis x=10,xmin=0.7,xmax=1400,xtick={1,10,100,1000},xticklabels={1,10,100,1000},ymin=-0.8,ymax=0.8,
  tick label style={font=\footnotesize},grid=major,grid style={gray!25},table/col sep=comma]
\addplot[omThm,thick,mark=*,mark size=1.3pt,error bars/.cd,y dir=both,y explicit] table[x=latency_ms,y=edge,y error=se]{figdata/hft/08-lead-lag-and-cross-venue-trading/edge.csv};
\draw[gray] (axis cs:0.7,0) -- (axis cs:1400,0);
\draw[omDef,dashed] (axis cs:500,-0.8) -- (axis cs:500,0.8);
\node[font=\scriptsize,omDef,anchor=north east,fill=white,inner sep=1pt] at (axis cs:480,0.75) {planted lead};
\end{axis}
\end{tikzpicture}

\omcaption{Edge per share of taking one lot of the follower after the leader has moved and the follower has not, after half the spread and a 0.1-tick
fee, marked five seconds later, against the taker's latency (log scale), $\pm$ one standard error; planted lead half a second, 135 trades in two
sessions of twenty minutes. Data: \texttt{hf\_leadlag.edge\_curve}.}\label{fig:hf:lead-lag-and-cross-venue-trading:edge}
\end{omfigure}

The edge is 0.485 ticks a share at one millisecond, 0.43 at 100, 0.24 at 300, and $-0.22$ at 500 milliseconds, the planted lead: the follower has
moved, the gap is gone, and the taker pays half the spread and the fee for nothing ($-0.54$ at one second). Most of the edge survives the first 100
milliseconds only because this market's lead is long. Budish, Cramton and Shim's numbers put the real future--fund lead at a few milliseconds in 2011,
so the same curve, drawn for a real pair, is compressed into the first milliseconds, and speed decides who earns it (chapter 9).

\section{When the leader changes}

A lead is a property of a pair at a time. The venue that leads can change with the hour (a futures market that opens before the stocks leads at the
open), with the move (large moves start in the deepest market), with a venue's fees or outages, and with competition: when enough traders trade the
lead, the follower's market makers withdraw their stale quotes faster (chapter 6), and the lead shortens. The practical consequences are to estimate
the lead continuously on recent data, by time of day, and to monitor the information share for a change of leader, as for any predictor's decay
(One Quant Book 7, chapter 13).

\section{Strategy files}

\begin{strategyfile}{Index future to fund}\label{strat:hf:lead-lag-and-cross-venue-trading:future}
\sfield{Who pays you, and why} The fund's liquidity providers whose quotes are stale for the milliseconds after the future moves.
\sfield{Instruments and venues} An index future and the exchange-traded fund on the same index (and its options).
\sfield{Signal} The future's move, in fund units through the fair-value basis (chapter 2), minus the fund's move over a short window.
\sfield{Sizing and execution} Take one clip of the fund when the gap exceeds half the spread plus fees; or, as the fund's market maker, withdraw the
stale side.
\sfield{Costs} Half the spread, taker fees, the hedge in the future if the position is held.
\sfield{How it dies} The arms race: the median opportunity lasted 97 milliseconds in 2005 and 7 in 2011, while its median profit per opportunity
stayed near 0.08 index points (Budish, Cramton and Shim); the edge goes to the fastest.
\sfield{Horizon, capacity, infrastructure} Milliseconds; a clip per opportunity; the fastest link between the future's and the fund's matching
engines (One Quant Book 14).
\sfield{Backtest honestly} Both venues' timestamps on one clock, the real propagation delay, and execution at the price available when the order
arrives.
\sfield{Sources} Hasbrouck (2003); Budish, Cramton and Shim (2015); this chapter: 0.485 ticks a share at one millisecond, gone before the planted lead.
\end{strategyfile}

\begin{strategyfile}{Fund to constituents}\label{strat:hf:lead-lag-and-cross-venue-trading:etf}
\sfield{Who pays you, and why} Market makers in single stocks whose quotes lag a market-wide move shown first in the fund or the future.
\sfield{Instruments and venues} A broad or sector fund and its largest constituents.
\sfield{Signal} The fund's move times each stock's beta, minus the stock's own move.
\sfield{Sizing and execution} Skew or withdraw the stale side of many stocks at once; take only in the largest gaps.
\sfield{Costs} Many small positions and their fees; a hedge in the fund.
\sfield{How it dies} Stock-specific news: a stock that did not follow may have its own reason. Sector funds follow the broad fund more than they lead
it (Hasbrouck), so the direction matters.
\sfield{Horizon, capacity, infrastructure} Milliseconds to seconds; a book of hundreds of names.
\sfield{Backtest honestly} Point-in-time betas; removing stocks with news in the window only if the live system could.
\sfield{Sources} Hasbrouck (2003).
\end{strategyfile}

\begin{strategyfile}{Cash Treasury to Treasury futures}\label{strat:hf:lead-lag-and-cross-venue-trading:treasury}
\sfield{Who pays you, and why} Liquidity providers in the slower of the cash note and its future, on whichever side leads at the time.
\sfield{Instruments and venues} On-the-run notes on the interdealer platforms (One Quant Book 2, chapter 4) and the Treasury futures.
\sfield{Signal} The leader's yield move translated by the conversion factor and the basis (One Quant Book 2, chapter 6) into the follower's price.
\sfield{Sizing and execution} Take in the follower; or skew quotes there.
\sfield{Costs} Platform fees, the basis's own moves, financing if held.
\sfield{How it dies} The leader changes: Mizrach and Neely found futures carrying most price discovery for five- and ten-year notes from 1999 on,
with the spot market's share varying with spreads, activity and volatility.
\sfield{Horizon, capacity, infrastructure} Milliseconds around data releases, seconds otherwise; access to both markets' matching engines.
\sfield{Backtest honestly} Estimate the leader on past data only, by period.
\sfield{Sources} Mizrach and Neely (2008).
\end{strategyfile}

\begin{strategyfile}{Leading crypto venue to lagging venue}\label{strat:hf:lead-lag-and-cross-venue-trading:crypto}
\sfield{Who pays you, and why} Makers on slower or smaller exchanges whose books follow the price-discovering venues late.
\sfield{Instruments and venues} One coin on several centralised exchanges (One Quant Book 3, chapter 16), spot and perpetual.
\sfield{Signal} The leading venue's move minus the follower's, estimated lead by venue pair.
\sfield{Sizing and execution} Take on the follower with inventory prepositioned there (One Quant Book 3, chapter 16); or make there with a skew.
\sfield{Costs} Taker fees in basis points, transfers to rebalance inventory, rate limits.
\sfield{How it dies} When capital moves freely the gaps close; Makarov and Schoar found deviations much larger across countries than within them, with
capital controls keeping them open.
\sfield{Horizon, capacity, infrastructure} Milliseconds to seconds; servers near each venue (One Quant Book 14).
\sfield{Backtest honestly} Each venue's timestamps and its data latency to the trader, not a merged tape.
\sfield{Sources} Makarov and Schoar (2020); chapter 24.
\end{strategyfile}

\begin{strategyfile}{Primary FX venue to secondary platforms}\label{strat:hf:lead-lag-and-cross-venue-trading:fx}
\sfield{Who pays you, and why} Liquidity providers on secondary platforms whose streams lag the primary venue (One Quant Book 2, chapter 14).
\sfield{Instruments and venues} Major currency pairs across the primary venue, other platforms and streams to aggregators.
\sfield{Signal} The primary venue's move minus the platform's, and triangular consistency across pairs.
\sfield{Sizing and execution} Take where the platform allows it; last look (chapter 9) lets the provider reject.
\sfield{Costs} Platform fees and rejections.
\sfield{How it dies} Algorithmic trading itself: Chaboud and co-authors found it reduced triangular arbitrage opportunities, mainly through computers
taking liquidity.
\sfield{Horizon, capacity, infrastructure} Milliseconds; the matching sites of the primary venue (One Quant Book 14, chapter 24).
\sfield{Backtest honestly} Include rejections under last look and hold times.
\sfield{Sources} Chaboud, Chiquoine, Hjalmarsson and Vega (2014).
\end{strategyfile}

\section{Tutorial: who moves first}\label{tut:hf:lead-lag-and-cross-venue-trading:tut}

\begin{tutorial}
\textbf{Goal.} Recover a planted lead from two asynchronous price series, measure price-discovery shares, and price the lead as a function of
latency. \textbf{End state:} \cref{fig:hf:lead-lag-and-cross-venue-trading:shares,fig:hf:lead-lag-and-cross-venue-trading:ccf,fig:hf:lead-lag-and-cross-venue-trading:edge}.
\begin{tutsteps}
\item \textbf{Markets.} \texttt{hf\_leadlag.pair(lag, seed)} simulates the leader and the follower, made responsive.
\item \textbf{Lead.} \texttt{estimate} applies Book 7's Hayashi--Yoshida estimator to the two mids' change points on a grid of candidate lags.
\item \textbf{Shares.} \texttt{firm.xvenue.shares} fits the error-correction model on a 0.1-second grid and returns the component share and the
information share's bounds (\cref{lst:hf:lead-lag-and-cross-venue-trading:shares}).
\omcode{code/firm/xvenue/firm_xvenue.py}{46}{71}{The error-correction model and the two price-discovery shares.}\label{lst:hf:lead-lag-and-cross-venue-trading:shares}
\item \textbf{Edge.} \texttt{lead\_edge} replays the two books without impact and trades the gap after a latency; \texttt{edge\_curve()} sweeps it.
\omcode{code/firm/xvenue/firm_xvenue.py}{81}{97}{Trade the gap after a latency, at the follower's price then; mark five seconds later.}\label{lst:hf:lead-lag-and-cross-venue-trading:edge}
\end{tutsteps}
\textbf{What to change next.} Make the lead random from move to move and see the cross-correlation's peak widen; give the follower its own
information (a share of its efficient price moves first) and watch the information shares meet.
\end{tutorial}

\section{Build: the cross-venue toolkit}\label{bld:hf:lead-lag-and-cross-venue-trading:xvenue}

\begin{build}
\bfield{Purpose} Who leads, by how much, and what the lead is worth at a given latency.
\bfield{Interface} \texttt{sample(t, x, grid)}; \texttt{lead(t1, x1, t2, x2, lags)}; \texttt{shares(p1, p2, lags)} with the component share and the
information-share bounds; \texttt{lead\_edge(lead\_t, lead\_mid, fol\_t, fol\_bid, fol\_ask, latency, move, window, H, fee)}.
\bfield{Rules} Leads are estimated on the series' own ticks; shares on a grid whose step is stated; edges after half the spread and fees, one lot,
no impact.
\bfield{Acceptance tests} \texttt{code/firm/xvenue/tests/}: on a random walk seen by one market three steps before the other, the leader's shares
exceed 0.8 and swap when the markets are swapped; a planted grid lead is recovered exactly; the edge of a hand-built pair at two latencies.
\bfield{Stretch} Shares by hour of the day; a lead that depends on the size of the move; several followers at once.
\end{build}

\begin{omsources}
\item J.~Hasbrouck, ``One security, many markets: determining the contributions to price discovery'', \emph{Journal of Finance} 50(4), 1995.
\item J.~Gonzalo and C.~Granger, ``Estimation of common long-memory components in cointegrated systems'', \emph{Journal of Business and Economic
Statistics} 13(1), 1995.
\item J.~Hasbrouck, ``Intraday price formation in U.S. equity index markets'', \emph{Journal of Finance} 58(6), 2003.
\item B.~Mizrach and C.~J.~Neely, ``Information shares in the US Treasury market'', \emph{Journal of Banking and Finance} 32(7), 2008.
\item I.~Makarov and A.~Schoar, ``Trading and arbitrage in cryptocurrency markets'', \emph{Journal of Financial Economics} 135(2), 2020.
\item A.~Chaboud, B.~Chiquoine, E.~Hjalmarsson and C.~Vega, ``Rise of the machines: algorithmic trading in the foreign exchange market'',
\emph{Journal of Finance} 69(5), 2014.
\item E.~Budish, P.~Cramton and J.~Shim, \emph{Quarterly Journal of Economics} 130(4), 2015.
\end{omsources}

\section{Exercises}

\begin{exercise}[$\star$]\label{exo:hf:lead-lag-and-cross-venue-trading:1}
An error-correction fit gives $\alpha_1=-0.02$ and $\alpha_2=0.18$. Which market leads, and what is its component share?
\end{exercise}

\begin{exercise}[$\star$]\label{exo:hf:lead-lag-and-cross-venue-trading:2}
A taker makes 135 trades in two twenty-minute sessions at 0.363 ticks a share, one lot of 100 shares, with a one-cent tick. What does it earn a
session?
\end{exercise}

\begin{exercise}[$\star$]\label{exo:hf:lead-lag-and-cross-venue-trading:3}
By what factor did the median life of the future--fund arbitrage shrink between 2005 and 2011?
\end{exercise}

\begin{exercise}[$\star\star$]\label{exo:hf:lead-lag-and-cross-venue-trading:4}
Why are the information share's bounds wide at a 50-millisecond lead and narrow at one second, on a 0.1-second grid?
\end{exercise}

\begin{exercise}[$\star\star$]\label{exo:hf:lead-lag-and-cross-venue-trading:5}
Why does the edge become negative, rather than zero, once the latency exceeds the lead?
\end{exercise}

\begin{exercise}[$\star\star$]\label{exo:hf:lead-lag-and-cross-venue-trading:6}
From \cref{fig:hf:lead-lag-and-cross-venue-trading:edge}, between which latencies does the edge cross zero?
\end{exercise}

\begin{exercise}[$\star\star\star$]\label{exo:hf:lead-lag-and-cross-venue-trading:7}
\emph{Coding.} With a planted lead of 0.2 seconds, compute the edge at latencies of one and of 100 milliseconds (\texttt{edge(0.2, 0.001)},
\texttt{edge(0.2, 0.1)}).
\end{exercise}

\begin{exercise}[$\star\star\star$]\label{exo:hf:lead-lag-and-cross-venue-trading:8}
\emph{Find the flaw.} ``On one-second data the future's information share is 99\%, so the fund contributes nothing to price discovery and we can
ignore its book.''
\end{exercise}

\section{Problem: Who Moves First}

\begin{problem}[{Weekend problem --- who moves first}]\label{pb:hf:lead-lag-and-cross-venue-trading:1}
A desk trades a fund against its index future and must decide how fast it needs to be.

\medskip\noindent\textbf{Part I --- The evidence.}
\begin{enumerate}
\item What did Hasbrouck find about price discovery in US equity index markets?
\item What did Mizrach and Neely find in Treasuries?
\item What happens to the future--fund correlation at 10 and at 1 millisecond?
\item Define a \omterm{def:hf:lead-lag-and-cross-venue-trading:lead}{cross-venue lead}.
\end{enumerate}

\medskip\noindent\textbf{Part II --- Measuring.}
\begin{enumerate}[resume]
\item Write the error-correction model and the two shares of \cref{prop:hf:lead-lag-and-cross-venue-trading:shares}.
\item Give the leader's shares at each planted lead and explain their pattern.
\item Why must the lead be estimated on the series' own ticks?
\item Give the lead estimates for each planted lead.
\item Why did the default simulated books hide a 50-millisecond lead?
\end{enumerate}

\medskip\noindent\textbf{Part III --- Trading.}
\begin{enumerate}[resume]
\item Describe the taker's rule and how its edge is measured.
\item Give the edge at 1, 100, 300 and 500 milliseconds and at one second.
\item Why is the edge still large at 100 milliseconds here, and why would it not be for a real future--fund pair?
\item How does a market maker in the follower use the same signal?
\end{enumerate}

\medskip\noindent\textbf{Part IV --- The verdict.}
\begin{enumerate}[resume]
\item State the \emph{named result}: the planted and estimated lead, the leader's information share, and the latency at which the edge disappears.
\item Why does the leader change over time?
\item Which strategy file depends most on capital controls, and why?
\item How would you monitor a change of leader in production?
\item Why does the edge not depend on the lead's length as long as the taker is faster than it?
\item What is the cost of a latency advantage that is not used?
\item In one sentence: what decides who earns a lead?
\end{enumerate}
\end{problem}

\section{Interview questions}

\begin{interviewq}[$\star$ \iqroles{trader}]\label{iq:hf:lead-lag-and-cross-venue-trading:1}
The future is up two ticks and the fund has not moved. List three reasons not to buy the fund.
\end{interviewq}

\begin{interviewq}[$\star\star$ \iqroles{researcher}]\label{iq:hf:lead-lag-and-cross-venue-trading:2}
Why does sampling two asynchronous series on a grid bias their correlation down, and how does the Hayashi--Yoshida estimator avoid it?
\end{interviewq}

\begin{interviewq}[$\star\star$ \iqroles{researcher}]\label{iq:hf:lead-lag-and-cross-venue-trading:3}
Explain why the information share has an upper and a lower bound, and what makes them close.
\end{interviewq}

\begin{interviewq}[$\star\star$ \iqroles{developer}]\label{iq:hf:lead-lag-and-cross-venue-trading:4}
Two venues' feeds arrive on servers in two data centres. How do you timestamp them so that a lead of a few milliseconds can be measured?
\end{interviewq}

\begin{interviewq}[$\star\star$ \iqroles{trader}]\label{iq:hf:lead-lag-and-cross-venue-trading:5}
Your edge from trading a lead halves in a month while its frequency is unchanged. What happened, and what do you check?
\end{interviewq}

\begin{interviewq}[$\star\star\star$ \iqroles{researcher}]\label{iq:hf:lead-lag-and-cross-venue-trading:6}
In a pair cointegrated with vector $(1,-1)$, show that if market 2 alone corrects ($\alpha_1=0$), market 1's component share is one, and discuss its
information share.
\end{interviewq}
